% C1_identification_v42.tex -- text v42 (fix22_texts_v42, 23 Nov 2026), built by fix22_texts_v42/tools/build_v42.py from fix21_texts_v41/C1_identification_v41.tex (unchanged):
% Lean replaces DKS (writeup/lean_bridge/REVIEW_W7.md section 3) and the v41 leftovers, where they touch this file.  Not compiled.
% C1_identification_v42.tex -- text v41 (fix21_texts_v41, 21 Nov 2026), built by fix21_texts_v41/tools/build_v41.py from fix19_texts_v40/C1_identification_v40.tex (unchanged):
% wording items W-1..W-5 and the retag item of the delta check fix20_delta_check/DELTA_CHECK_v40.md where they touch this file.  Not compiled.
% C1_identification_v42.tex -- text v40 (fix19_texts_v40, 21 Nov 2026), built by fix19_texts_v40/tools/build_v40.py from fix17_texts_v39/C1_identification_v39.tex (unchanged):
% fixes of the final review fix18_final_review/FINAL_REVIEW_v39.md (F-1..F-7, R-1..R-3) where they touch this file.  Not compiled.
% C1_identification_v42.tex -- text v39 (fix17_texts_v39, 20 Nov 2026), built by fix17_texts_v39/tools/build_v39.py from fix16_texts_v38/C1_identification_v38.tex (unchanged):
% re-pinned step-4 records (manifest_v38 be909bbf...), v4 kits (WSL test runs only), step4_handoff_v42.tex; v38 retag fixes.  Not compiled.
% C1_identification_v42.tex -- text v38 (fix16_texts_v38), 20 Nov 2026.  Built by tools/build_interface_v38.py from
% fix14_interface/C1_identification.tex (fix round 14, unchanged).  New in v38: Lemma (P2-a) with proof (review fix15 C1-1), the
% Y4 label (C1-2), the x/y wording of Lemma kP2 (C1-3).  Original header follows.
% (fix round 14)  Closes review item C-1 of fix11_review/REVIEW_v4b_v36_proofs.md:
% smoothness of gamma(t_1) and the identification of step 4's validated solution with the smooth CDF flow.  New file only; nothing closed
% is edited.  Not compiled (no LaTeX in WSL).  Checks: C1_sym.py -> C1_sym.json (sympy, all_pass), C1_ke.py -> C1_ke.json (arb),
% C1_num.py -> C1_num.json (float sanity check only).
\documentclass[11pt]{article}
\usepackage[margin=1in]{geometry}
\usepackage{amsmath,amssymb,amsthm,enumitem}
\newtheorem{lemma}{Lemma}[section]\newtheorem{proposition}[lemma]{Proposition}\newtheorem{theorem}[lemma]{Theorem}\newtheorem{remark}[lemma]{Remark}
\newcommand{\fn}[1]{\texttt{\detokenize{#1}}}
\newcommand{\tH}{\textup{[hand]}}\newcommand{\tS}{\textup{[sym]}}\newcommand{\tC}{\textup{[cert]}}\newcommand{\tI}{\textup{[cited]}}\newcommand{\tN}{\textup{[num]}}\newcommand{\tL}{\textup{[Lean]}}
\providecommand{\LeanSthree}{\textup{[Lean]}}% S3 items of review W7 (lean_bridge snap_s3); text v42 first marked them [hand; Lean formalization in progress]
\newcommand{\ip}[2]{\langle#1,#2\rangle}
\begin{document}
\title{C-1: the step-4 solution is the smooth curve diffusion flow from $\Gamma_0$,\\ and $\gamma(t_1)$ is a smooth symmetric immersion}
\date{23 November 2026 (text v42; text v41; text v40; text v39; text v38 added (P2-a); first version fix round 14). \emph{Not compiled.}}
\maketitle

\noindent\textbf{Labels.} \tS: exact sympy identity in \fn{C1_sym.py} (\fn{C1_sym.json}, \fn{all_pass: true}; items Y1--Y10). \tC: ball arithmetic (\fn{python-flint arb}), script and JSON named. \tI: a result of another document, used as stated there (with its own label). \tL\ (text v42): a compiled theorem of Micah Warren's Lean development, as listed in \fn{U1_regularity_v42.tex} \S1.2 (directory \fn{CWFlow/Curve/}, commit pending, files identified by the SHA256 sums of \fn{writeup/lean_bridge/snap_final/SHA256SUMS}; compiled per file; axioms \fn{propext}, \fn{Classical.choice}, \fn{Quot.sound}); text v41 cited the Lean theorems as \tI, ``statement read from source, not compiled''. \LeanSthree: the S3 items of review W7 (S3a: periodic curve to \fn{Circle}; S3b: the shifted maximal flow and endpoint gluing), formalized after that review in \fn{Periodic2Circle.lean} and \fn{MaxShift.lean} (snapshot \fn{writeup/lean_bridge/snap_s3/SHA256SUMS}), printed as [Lean]; one macro, kept so that the provenance stays visible. \tH: hand argument written here, not machine-checked. \tN: float sanity check, not a proof step.

\noindent\textbf{What is closed.} Step~4 (\fn{step4_proof.tex} l.130, l.336, l.370, l.672) asserts the smoothness of $\gamma$ for $t>0$ ``by parabolic regularity'' and the continuation of the P2 solution by ``the literature fact (P1)'', both labelled [A]; and no text identifies its validated object (the P1 Galerkin solution in $C^0H^5\cap C^1L^2$, in its own gauge, then the P2 normal graph over the base) with a smooth flow. Theorem~\ref{C1:thm} below does this, with existence, uniqueness, maximality, symmetry and the blow-up criterion from the compiled Lean theorems (T1), (Sy) of \fn{U1_regularity_v42.tex} \S1.2 (text v42; text v41: (L1), (L2) read from source and DKS (D1)), and step~4's certified bounds. \emph{New in text v38:} the a-priori form of Theorem~P2 is Lemma~\ref{C1:lem:P2a} with a written proof (\S\ref{C1:sec:P2a}; review fix15 item C1-1); in fix round~14 it was a reading of step~4's Step~7. It also gives (S4.1) and (S4.2) of \fn{fix22_texts_v42/U1_regularity_v42.tex}, and it shows that the physical solution of step~3 is the continuation of the same flow.

%======================================================================
\section{Inputs}\label{C1:sec:in}
\paragraph{(I1) $\Gamma_0$ \tI.} Theorem~4 and Proposition ``$\Gamma_0$'' of \fn{step4_proof.tex} (\S``The initial curve''; certificate \fn{t2/gamma0_final.json}): $\theta_0(\alpha)=-\frac\pi2+\alpha+\mathrm{th}_0(\alpha)$, $\mathrm{th}_0=\sum_{n\le3762}b_n\sin n\alpha$ with exact doubles $b_n$ ($n\ne2$) and $b_2=b_2^*$ the certified unique zero of $C(b_2)=\int_0^\pi\cos(\theta_0+\frac\pi2)\,d\alpha$; $L_0$, $X_0(0)$ real. Put $X^0(\alpha)=X_0(0)+\frac{L_0}{2\pi}\int_0^\alpha e^{i\theta_0}$.
\begin{lemma}\label{C1:lem:G0}
$X^0$ is $2\pi$-periodic, real-analytic, $|X^0_\alpha|=L_0/2\pi$, and $X^0(-\alpha)=\overline{X^0(\alpha)}$. So $F_0:\mathbb S^1\to\mathbb C$, $F_0(e^{i\alpha}):=X^0(\alpha)$, is a smooth immersion with $\bar F_0=F_0\circ\iota$, $\iota(z)=\bar z$.
\end{lemma}
\begin{proof}\tH, \tS\ Y10. Since $\mathrm{th}_0$ is odd, $\theta_0(-\alpha)=-\pi-\theta_0(\alpha)$ (Y10a), so $e^{i\theta_0(-\alpha)}=-e^{-i\theta_0(\alpha)}$ (Y10b), and $X^0(-\alpha)=X_0(0)+\frac{L_0}{2\pi}\int_0^\alpha e^{-i\theta_0}=\overline{X^0(\alpha)}$. Periodicity: $\int_0^{2\pi}e^{i\theta_0}=\int_{-\pi}^{\pi}e^{i\theta_0}=\int_0^\pi(e^{i\theta_0}-e^{-i\theta_0})=2i\int_0^\pi\sin\theta_0=-2i\,C(b_2^*)=0$. $\theta_0-\alpha$ is a trigonometric polynomial, so $X^0$ is real-analytic; $|X^0_\alpha|=L_0/2\pi>0$.
\end{proof}

\paragraph{(I2) P1 and the Galerkin-limit lemma \tI.} \fn{step4_proof.tex} \S``P1'' and the appendix ``Proof of the Galerkin-limit lemma'' (ll.~236--247, 822--831). In the sine class $\mathrm{th}=\sum_{n\ge1}u_n\sin n\alpha$, with clock $d\sigma=(2\pi/L)^4dt$,
\begin{equation}\label{C1:eq:P1}
\mathrm{th}_\sigma=-\mathrm{th}''''-\mathrm{th}''+N(\mathrm{th}),\qquad N=-\bigl(v^2+\tfrac13v^3\bigr)'+(1+v)\,P\bigl(w^2-\ip{w^2}{}\bigr),\quad v=\mathrm{th}',\ w=\mathrm{th}'',
\end{equation}
$P$ the zero-mean antiderivative, $\ip{\cdot}{}$ the mean; slaved: $(\ln L)'=-\frac12\sum n^4u_n^2$, $t'=(L/2\pi)^4$, $X_0'=(L/2\pi)\sum n^3u_n$. The 96 slabs are cut into finitely many sub-intervals $I_k$ ($k=1,\dots,K$; ten per slab). On each, the lemma constructs a solution $u\in\mathcal C(I_k):=C^0(I_k;H^5)\cap C^1(I_k;L^2)$ of \eqref{C1:eq:P1} in the sine class, starting from the end value of $I_{k-1}$ (from $\mathrm{th}_0$ for $k=1$). The concatenation $u$ on $[0,\sigma_{96}]$ and the slaved $L,t,X_0$ define the \emph{P1 curve} $X^{\rm P1}(\sigma,\alpha)=X_0(\sigma)+\frac{L(\sigma)}{2\pi}\int_0^\alpha e^{i(-\pi/2+\alpha'+u)}d\alpha'$; $t(\sigma_{96})=t_{96}$, and $L(\sigma_{96})=48.293\ldots>0$.

\paragraph{(I3) Theorem P2 \tI.} \fn{step4_proof.tex} \S``P2'' (Theorem P2 and its proof, Steps 1--7; status there: constants [C], analytic steps [D], written out in \fn{LEMMA_F.md}, \fn{T7_REPORT.md}). Read precisely, Steps 1--6 and the continuity part of Step~7 are a priori statements about a given smooth CDF solution: Step~1 writes it as a normal graph $e$ over the base, Step~2 obtains $(\lambda,c,t)$ as $C^1$ functions of the label $\tau$ from the implicit function theorem (valid while $\bar m\|E\|<\frac12$; at the box $\bar m\|\hat E\|\le1.04\cdot10^{-4}$), Steps 3--6 bound the energies, and the continuity argument shows that the boxes $\mathbb B\cup\mathbb B'$ improve strictly. The only non-a-priori step is the last sentence of Step~7 (``the solution continues: this is the continuation fact (P1) of step~3, from the literature [A]''). We use P2 in the a-priori form of Lemma~\ref{C1:lem:P2a} (\S\ref{C1:sec:P2a}), which writes out the continuity argument of Step~7: the limits of $(\lambda,c,t)$ at the end of the interval, closedness and openness of the global normal-graph representation over the base, and the implicit function theorem at the limit.

\paragraph{(I4) Lean \tL.} (Text v42; review \fn{writeup/lean_bridge/REVIEW_W7.md}: sound, no blocking finding.) The theorems (L1) \fn{cdf_exists}, (L2) \fn{cdf_unique}, (T1) \fn{cdf_exists_maximal}, (T2) \fn{cdf_extends_of_bounded_energy}, (V) \fn{flowVel_eq_neg_kss}, (Sy) \fn{cdf_symmetric}, (Sh) \fn{IsFlow.shift_gen}, (PC) \fn{exists_isLagImmersion_of_periodic}, (MS) \fn{cdf_maximal_shift}, (GE) \fn{isFlow_glue_endpoint} of \fn{CWFlow/Curve/} and (Gl) \fn{rst_glue} of the library, as stated in \fn{U1_regularity_v42.tex} \S1.2: $M=\mathbb S^1$ is Mathlib's \fn{Circle}, $V=\mathbb C$, $J=i$. They are compiled per file at $n=1$ (commit pending, SHA256 sums in \fn{writeup/lean_bridge/snap_final/SHA256SUMS} and \fn{writeup/lean_bridge/snap_s3/SHA256SUMS}; a full \fn{lake build CWFlow.Curve.All} is pending). Text v41's caveat ``read from source, not compiled; the $n=1$ instance not elaborated'' (review C-3) no longer applies. \fn{IsFlow E1 Jc T Phi}: $\Phi$ jointly $C^\infty$ on $[0,T]\times\mathbb S^1$, immersions, $\partial_t\Phi=\fn{flowVel}(\Phi)$; at $n=1$, $\fn{flowVel}(X)=-k_{ss}N$, $N=iT$, with no tangential part ((V); Lemma~``$n=1$'' of \fn{U1_regularity_v42.tex}, with S1--S3 of \fn{U1_sym.py} \tS\ as a cross-check). The maximal flow, its uniqueness, its symmetry and the blow-up criterion are (T1) and (Sy) for any $F_0$ with \fn{IsLagImmersion E1 Jc F0}; no step-3 statement is involved.

\paragraph{(I5) Criterion \tL.} (Text v42.) The blow-up criterion is (T1)(iii): if $T_*<\infty$, then $\fn{bendingEnergy}(\fn{lift}(\gamma(t)))=\int_{\gamma(t)}k^2ds$ is unbounded on $[0,T_*)$. Text v41 used here DKS (D1), a step inside the proof of DKS Theorem~3.1, with a hand completion (review item C-4: the extension is an immersion at the extension time). Neither is used now. If an extension form is wanted, it is (T2), whose extension is an \fn{IsFlow} on a closed interval $[0,T+\varepsilon]$, hence an immersion at $T$ with the velocity identity there \tL. C-1 uses only (T1), through Lemma~\ref{C1:lem:max}.

%======================================================================
\section{The smooth flow from $\Gamma_0$}
\begin{lemma}[maximal flow]\label{C1:lem:max}
There are $T_*\in(0,\infty]$ and a unique $\gamma:[0,T_*)\times\mathbb S^1\to\mathbb C$ with $\gamma|_{[0,T]}$ an \fn{IsFlow} from $F_0$ for every $T<T_*$, and no \fn{IsFlow} from $F_0$ on $[0,T_*]$ if $T_*<\infty$. It satisfies $\overline{\gamma(t)}=\gamma(t)\circ\iota$, and if $T_*<\infty$ then $\sup_{t<T_*}\int_{\gamma(t)}k^2ds=\infty$.
\end{lemma}
\begin{proof}\tL\ (text v42). \emph{The data on \fn{Circle}} \LeanSthree\ (review W7 S3a): $X^0$ is $C^\infty$ and $2\pi$-periodic with $X^0_\alpha\ne0$ (Lemma~\ref{C1:lem:G0}), so (PC) \fn{exists_isLagImmersion_of_periodic} gives $F$ with $\fn{lift}\,F=X^0$ and \fn{IsLagImmersion E1 Jc F}; $\fn{lift}\,F=X^0$ says $F(e^{i\alpha})=X^0(\alpha)$, so $F=F_0$ (\fn{lift_injective}). \emph{The rest} \tL: (T1) with this $F_0$ gives $T_*$ and $\gamma$; clauses (i), (ii) give existence on every $[0,T]$, $T<T_*$, and that every \fn{IsFlow} from $F_0$ on $[0,T]$ has $T<T_*$ and is the restriction of $\gamma$ (so none exists on $[0,T_*]$, and $\gamma$ is unique on $[0,T_*)$); clause (iii) is the criterion, with $\fn{bendingEnergy}(\fn{lift}(\gamma(t)))=\int_{\gamma(t)}k^2ds$. Symmetry: $\bar F_0=F_0\circ\iota$ (Lemma~\ref{C1:lem:G0}) is the hypothesis of (Sy), and (Sy) on each $[0,T]$, $T<T_*$, gives $\overline{\gamma(t)}=\gamma(t)\circ\iota$. (Text v41: the U-1 lemmas assembled by hand from (L1), (L2), gluing, and DKS (D1) with the C-4 completion.)
\end{proof}

\noindent\emph{Gauge.} For $t<T_*$ let $p_0=1\in\mathbb S^1$ (a fixed point of $\iota$; $\alpha=0$ of $\Gamma_0$), $s(t,p)=\int_{p_0}^p|\gamma_p|$ (counter-clockwise), $L(t)=s(t,p_0^{+2\pi})$, $\alpha(t,p)=2\pi s/L$, $\theta(t,\alpha)$ the continuous lift of the tangent angle of $\gamma(t)$ at the point with normalised arclength $\alpha$, normalised by continuity in $t$ from $\theta(0,\cdot)=\theta_0$, and $\mathrm{th}^\gamma(t,\alpha)=\theta(t,\alpha)+\frac\pi2-\alpha$; $X^\gamma_0(t)=\gamma(t,p_0)$; $\sigma^\gamma(t)=\int_0^t(2\pi/L)^4$.

\begin{lemma}[the flow in step 4's gauge]\label{C1:lem:gauge}
For every $T<T_*$: (a) $\alpha$ is jointly $C^\infty$ with $\partial_p\alpha>0$, so $(t,\alpha)\mapsto\mathrm{th}^\gamma$ is jointly $C^\infty$ on $[0,T]\times\mathbb S^1$; (b) $\mathrm{th}^\gamma(t,\cdot)$ is odd and $2\pi$-periodic, i.e.\ in the sine class, and $\theta(t,0)=-\frac\pi2$; (c) as a function of $\sigma=\sigma^\gamma(t)$, $\mathrm{th}^\gamma$ solves \eqref{C1:eq:P1} pointwise, and $L$, $t$, $X_0^\gamma$ solve the three slaved equations of (I2); (d) $\mathrm{th}^\gamma\in\mathcal C([0,\sigma^\gamma(T)])$.
\end{lemma}
\begin{proof}
(a) \tH\ $|\gamma_p|$ is smooth and positive, so $s,L,\alpha$ are smooth, $L>0$, $\partial_p\alpha=2\pi|\gamma_p|/L>0$; the inverse $p(t,\alpha)$ is smooth by the inverse function theorem, and the tangent angle is a smooth function of $\gamma_p$ composed with $p(t,\alpha)$.

(b) \tH, \tS\ Y10. Differentiate $\overline{\gamma(t)}=\gamma(t)\circ\iota$ (Lemma~\ref{C1:lem:max}): $|\gamma_p|$ is $\iota$-invariant, so $s(t,\iota p)=-s(t,p)$ (mod $L$) and $\alpha\circ\iota=-\alpha$; the tangents satisfy $\overline{T(\alpha)}=-T(-\alpha)$, i.e.\ $\theta(t,-\alpha)\equiv-\pi-\theta(t,\alpha)\pmod{2\pi}$. At $\alpha=0$ this gives $\theta(t,0)\in-\frac\pi2+\pi\mathbb Z$; it is continuous in $t$ and equals $-\frac\pi2$ at $t=0$, so $\theta(t,0)=-\frac\pi2$. Both sides of $\theta(t,-\alpha)=-\pi-\theta(t,\alpha)$ are continuous in $\alpha$ and agree at $\alpha=0$, so they agree identically: $\mathrm{th}^\gamma$ is odd (Y10a). The rotation index $(\theta(t,2\pi)-\theta(t,0))/2\pi$ is an integer depending continuously on $t$, equal to $1$ at $t=0$; so $\mathrm{th}^\gamma$ is periodic.

(c) \tS\ Y1--Y7, \tH\ for the bookkeeping; for Y4 the inputs $\partial_ts|_p=-\int_0^skV$ and $L_t=-\int kV$ are typed into \fn{C1_sym.py} ll.72--76, so the drift formula is \tH\ (inputs, derived from Y1b two lines above) plus \tS\ (the algebra); review fix15 C1-2. With $V=-k_{ss}$ the normal speed, Y1a gives $\theta_t|_p=\partial_sV=-k_{sss}$ and Y1b gives $\partial_t|\gamma_p|=-kV|\gamma_p|$; hence $\partial_ts|_p=-\int_0^skV\,ds'$, $L_t=-\int kV\,ds$. In the gauge, Y2 expresses $k,k_s,k_{ss},k_{sss}$ through $v,w$; Y4 gives $\partial_\sigma\alpha|_p=\int_0^\alpha(1+v)w'+\alpha\ip{w^2}{}$ (with Y4b for the integrand and Y3 over a period for $L_t$); then $\theta_\sigma|_\alpha=\theta_\sigma|_p-(\partial_\sigma\alpha|_p)\theta_\alpha$ (chain rule) and Y3 (integrated from $0$, with $w(0)=0$ by oddness, Y6e) give Y5a:
\[\mathrm{th}^\gamma_\sigma=-\mathrm{th}''''-(1+v)^2w+(1+v)\int_0^\alpha\bigl(w^2-\ip{w^2}{}\bigr).\]
Y5b rewrites $(1+v)^2w=\mathrm{th}''+(v^2+v^3/3)'$. For odd $\mathrm{th}$, $w^2-\ip{w^2}{}$ is even and $\int_0^\alpha(w^2-\ip{w^2}{})$ is odd with mean $0$, hence equals $P(w^2-\ip{w^2}{})$ (Y6b--d, checked on a symbolic odd trigonometric polynomial of degree 3; the general case is the same two-line parity argument \tH). This is \eqref{C1:eq:P1}. Slaved equations: $\frac{d}{d\sigma}\ln L=-\frac1{2\pi}\int w^2=-\frac12\sum n^4u_n^2$ (Y7a, Y6a; the Parseval step for infinitely many modes is \tH); $X_0^\gamma$ is the image of the fixed parameter $p_0$, so $\partial_tX^\gamma_0=V(p_0)N(p_0)$ with $N(p_0)=ie^{-i\pi/2}=1$ (Y7c) and $\frac{d}{d\sigma}X^\gamma_0=-(L/2\pi)\mathrm{th}'''(0)=(L/2\pi)\sum n^3u_n$ (Y7b); $t'=(L/2\pi)^4$ by definition of $\sigma$.

(d) \tH\ $\sigma^\gamma$ is a $C^1$ diffeomorphism of $[0,T]$ onto $[0,\sigma^\gamma(T)]$, and a jointly smooth function on a compact set is in $C^0H^5\cap C^1L^2$.
\end{proof}
\noindent\tN\ \fn{C1_num.py}: one explicit normal step of a closed symmetric test curve, re-gauged numerically, reproduces \eqref{C1:eq:P1} on the low modes with relative defect $3.0\cdot10^{-6}$ ($dt=10^{-6}$) and $3.1\cdot10^{-7}$ ($dt=10^{-7}$), and $(\ln L)'$ with $3.6\cdot10^{-5}$ ($dt=10^{-7}$). This is independent of the derivation above (it never uses Y4/Y5) and is only a cross-check.

\begin{lemma}[uniqueness in $\mathcal C$, without a box]\label{C1:lem:uniq}
Let $u,\tilde u\in C^0([0,S];H^5)\cap C^1([0,S];L^2)$ be odd, solve \eqref{C1:eq:P1} in $L^2$ for every $\sigma$, and $u(0)=\tilde u(0)$. Then $u\equiv\tilde u$.
\end{lemma}
\begin{proof}\tH\ (the argument of step~6 of the Galerkin-limit proof, \fn{step4_proof.tex} l.829, written for the whole class: the lemma's statement claims uniqueness only among solutions staying in its isolating set $W$, while its proof does not use $W$.) Let $\delta=u-\tilde u$, $f=\|\delta\|^2$ ($L^2(0,2\pi)$), so $f\in C^1$, $f'=2\ip{\delta_\sigma}{\delta}$. Let $M:=\sup_{[0,S]}(\|v\|_\infty+\|\tilde v\|_\infty+\|w\|+\|\tilde w\|)<\infty$ ($H^5\subset C^4$, continuity on a compact interval). Periodic integration by parts: $\ip{-(\partial^4+\partial^2)\delta}{\delta}=-\|\delta''\|^2+\|\delta'\|^2$. With $g(v)=v^2+v^3/3$, $|g(v)-g(\tilde v)|\le(2M+M^2)|\delta'|$, so $|\ip{-(g(v)-g(\tilde v))'}{\delta}|=|\ip{g(v)-g(\tilde v)}{\delta'}|\le(2M+M^2)\|\delta'\|^2$. Next $(1+v)P(w^2)-(1+\tilde v)P(\tilde w^2)=\delta'P(w^2)+(1+\tilde v)P(\delta''(w+\tilde w))$ ($P$ kills constants), $\|Pf\|_\infty\le4\|f\|_{L^1}$, so this pairs with $\delta$ to at most $4\sqrt{2\pi}M^2\|\delta'\|\|\delta\|+8\sqrt{2\pi}(1+M)M\|\delta''\|\|\delta\|$. With $\|\delta'\|^2\le\|\delta\|\|\delta''\|$ and Young, $\frac12f'\le-\frac12\|\delta''\|^2+C(M)f$; Gr\"onwall and $f(0)=0$ give $f\equiv0$.
\end{proof}

%======================================================================
\section{Identification on P1 ($0\le t\le t_{96}$)}
\begin{lemma}[curvature bound on P1]\label{C1:lem:kP1}
$\sup_{\sigma\le\sigma_{96}}\int_{X^{\rm P1}(\sigma)}k^2ds<\infty$.
\end{lemma}
\begin{proof}\tH\ $\int k^2ds=\frac{2\pi}L\int_0^{2\pi}(1+v)^2d\alpha$ (Y2a), $L$ is nonincreasing (slaved equation) so $L\ge L(\sigma_{96})>0$, and $\sup_{[0,\sigma_{96}]}\|v\|_\infty<\infty$ because $u\in C^0(I_k;H^5)$ on each of the finitely many $I_k$.
\end{proof}

\begin{proposition}[identification on P1]\label{C1:prop:P1}
$T_*>t_{96}$, $\sigma^\gamma(t_{96})=\sigma_{96}$, and for $t\in[0,t_{96}]$: $\mathrm{th}^\gamma(t)=u(\sigma^\gamma(t))$, $L$, $X_0$ agree, and $\gamma(t)$ is the P1 curve up to the reparametrisation $p\mapsto\alpha(t,p)$. In particular the P1 enclosures are enclosures of the smooth flow.
\end{proposition}
\begin{proof}\tH\ For $t'<T_*$ let $S'=\min(\sigma^\gamma(t'),\sigma_{96})$. By Lemma~\ref{C1:lem:gauge}, $\mathrm{th}^\gamma$ is in the class of Lemma~\ref{C1:lem:uniq} on $[0,S']$ and starts from $\mathrm{th}_0$; applying Lemma~\ref{C1:lem:uniq} on $I_1\cap[0,S'],I_2\cap[0,S'],\dots$ in turn gives $\mathrm{th}^\gamma=u$ on $[0,S']$. Then $L,t,X_0$ agree on $[0,S']$ (same ODEs, same data), so the curves agree. Case A: $\sigma^\gamma(t')>\sigma_{96}$ for some $t'<T_*$; then the flow time at $\sigma_{96}$ is $t(\sigma_{96})=t_{96}<t'<T_*$, which is the claim. Case B: $\sigma^\gamma(t)\le\sigma_{96}$ for all $t<T_*$; then every $\gamma(t)$, $t<T_*$, is a P1 curve, so $\int k^2ds$ is bounded (Lemma~\ref{C1:lem:kP1}) and $T_*=\infty$ (Lemma~\ref{C1:lem:max}); but $L\le L_0$ gives $\sigma^\gamma(t)\ge(2\pi/L_0)^4t\to\infty$, a contradiction.
\end{proof}

%======================================================================
\section{Identification on P2 ($t_{96}\le t\le t_1$)}
\paragraph{What P2 certifies.} Theorem~P2 is a statement about a smooth (classical) CDF solution, used only through geometric quantities: its normal velocity (Step~1 uses $\ip{\Gamma_\tau}{\nu}$ only), its normal graph $e$ over the base, and energies of $e$ (Step~1: $e\in W^{7,\infty}$ in $x$, identities a.e.\ in $\tau$). These are independent of the parametrisation, so the Lean flow, whose velocity is purely normal, is directly in P2's class; no gauge change is needed (unlike P1). Its initial datum is the P1 state at $t_{96}$ (switch lemma), which is $\gamma(t_{96})$ by Proposition~\ref{C1:prop:P1}.

\begin{lemma}[the normal-graph factor]\label{C1:lem:ke}
On the boxes $\mathbb B\cup\mathbb B'$ of P2, for every slab $2\ldots180$ and every stored cell of the corrected base, $|\tilde\kappa\,e|\le1.743\cdot10^{-3}$; hence $a:=1-\tilde\kappa e\ge\frac12$ and $J:=(a^2+e'^2)^{1/2}\ge\frac12$.
\end{lemma}
\begin{proof}\tC\ \fn{C1_ke.py} $\to$ \fn{C1_ke.json} (\fn{sup_ke} $\le1.7429\cdot10^{-3}$ at slab 180 (text v38 printed $=1.7428\cdot10^{-3}$, the 5-digit truncation; pinned as stage 22 of \fn{FINAL_manifest_v38.json} (29 stages, \fn{be909bbf...}), \fn{out_v38/iface/C1_ke.json}, equal to this file apart from args), $x\approx3.6\cdot10^5$; box $|e|\le32.3$ there): hulls of $\tilde\kappa$ from \fn{p2/t6c/jets} (widened outward by 2 ulps), $|e|\le10^{-5}$ for $y<28$, $|e|\le1.7E(x)$ for $y\ge28$ with $\sqrt{F_0}\le a_F\bar Y+b_F$ per slab (\fn{amp_A1_w180.json}) at $\bar Y=1.7\cdot10^{-12}$ (conservative). All 180 input files are hashed in the JSON.
\end{proof}


\subsection{(P2-a): Theorem P2 as an a-priori statement}\label{C1:sec:P2a}
\noindent\emph{Notation.} For $p=(\lambda,c,t)$ with $\lambda>0$ and $t\in[t_{96},T')$ put $\Gamma_p:=(\gamma(t)-ce_1)/\lambda$. For a label $\tau\in[\tau_a,27.15]$ let $\mathcal G(\tau,p)$ be the statement: there is $e\in C^1([0,\tilde\ell(\tau)])$ with $\sup|\tilde\kappa e|\le\frac12$ such that $x\mapsto\tilde\Gamma(x,\tau)+e(x)\tilde\nu(x,\tau)$, continued by the mirror symmetry, is a $C^1$ reparametrisation of the closed curve $\Gamma_p$ (a \emph{global} normal graph over the closed base). Under $\mathcal G$ put $F_j(\tau,p):=\ip{\chi e}{\varphi_j}_\rho$, $j=1,2,3$. The boxes $\mathbb B\cup\mathbb B'$ (with $|\vartheta|\le10^{-3}$) and the energies $\Lambda_e$, $F_j$ are those of \fn{step4_proof.tex} \S``Objects''. The base is $C^4$ in $x$ with enclosed jets that are continuous in $\tau$, and $C^1$ in $\tau$ on each slab (\fn{step4_proof.tex} \S``The base'': Radau collocation per slab, continuous across slabs).
\begin{lemma}[(P2-a)]\label{C1:lem:P2a}
Let $T'\in(t_{96},\infty]$ and let $\gamma$ be a smooth CDF flow on $[t_{96},T')$ (jointly $C^\infty$, immersed, $\overline{\gamma(t)}=\gamma(t)\circ\iota$) with $\gamma(t_{96})$ the P1 curve at $\sigma_{96}$. Then there are $\tau_e\in(\tau_a,27.15]$ and continuous functions $\lambda>0$, $c$, $t$ on $[\tau_a,\tau_e)$, $C^1$ on each base slab, with $t(\tau_a)=t_{96}$, $t(\tau)<T'$, $t_\tau=\lambda^4(1-\vartheta)$ and $|\vartheta|\le3.37\cdot10^{-8}$, such that on $[\tau_a,\tau_e)$: $\mathcal G(\tau,p(\tau))$ holds, with $e\in W^{7,\infty}$ in $x$; the constraints $F_j(\tau,p(\tau))=0$ hold; $\sqrt{\Lambda_e}\le\bar Y=10^{-12}$ and $\sqrt{F_j}\le A_j\bar Y+B_j$; and the boxes $\mathbb B\cup\mathbb B'$ hold. Moreover exactly one of the following holds:
\begin{enumerate}[label=(\roman*),leftmargin=2em]
\item $\tau_e=27.15$, $t(27.15)<T'$, and all of the above holds on $[\tau_a,27.15]$;
\item $t(\tau)\uparrow T'$ as $\tau\uparrow\tau_e$.
\end{enumerate}
\end{lemma}
\begin{proof}
\emph{Inputs and labels.} Steps 2--6 of the proof of Theorem~P2 and the pointwise bounds of its \S``boxes'' are used as stated in step~4 (constants \textup{[C]}, analytic steps \textup{[D]}, written out in \fn{LEMMA_F.md}, \fn{T7_REPORT.md}) \tI; the switch lemma (\fn{step4_proof.tex} \S``Switch'') gives the decomposition at $\tau_a$ \tI; Lemma~\ref{C1:lem:ke} is \tC. Everything else below (Steps a--f) is \tH, written here and not machine-checked.

\emph{(a) Local structure.} If $\mathcal G(\tau_0,p_0)$ holds with $\sup|\tilde\kappa e|\le c_1<1$, then the normal coordinates $(x,n)\mapsto\tilde\Gamma(x,\tau_0)+n\tilde\nu(x,\tau_0)$ of the closed base are a local diffeomorphism near the graph, and the foot-point map $\theta\mapsto X(\tau,p,\theta)$ from the parameter circle of $\Gamma_p$ to the base satisfies $\partial_\theta X=|\partial_\theta\Gamma_p|\cos d/(1-\tilde\kappa n)>0$ at $(\tau_0,p_0)$ ($d$ the angle between the tangents; $|\tan d|=|e'|/(1-\tilde\kappa e)$ is finite). The base is a closed curve, so there is no endpoint; by compactness and the implicit function theorem (as in Lemma GS-3 of \fn{GS_lemma_v42.tex}, here without the endpoint step) $\mathcal G(\tau,p)$ holds for $(\tau,p)$ near $(\tau_0,p_0)$, the foot-point map stays a degree-one diffeomorphism, and $e=\mathcal E(\tau,p;x)$ is $C^\infty$ in $p$, $C^4$ in $x$ with $x$-jets up to order $4$ continuous in $(\tau,p)$, and $C^1$ in $\tau$ on each slab. Hence $F=(F_1,F_2,F_3)$ is continuous in $\tau$ and $C^\infty$ in $p$, and $D_pF$ is the matrix $M+E$ of Step~2 (up to the fixed scaling of the directions $g_i$), invertible while $\bar m\|E\|<\frac12$.

\emph{(b) The set.} Let $\mathcal S$ be the set of $\tau'\in[\tau_a,27.15]$ for which there is a continuous $p$ on $[\tau_a,\tau']$, $C^1$ on each slab, with $p(\tau_a)=p_a$ (the parameters of the switch lemma), $t(\tau)<T'$, $\mathcal G(\tau,p(\tau))$, $F(\tau,p(\tau))=0$ and $\mathbb B\cup\mathbb B'$ for $\tau\in[\tau_a,\tau']$. Such a $p$ is unique: for two of them the coincidence set contains $\tau_a$, is closed, and is open by the local uniqueness in the implicit function theorem of (a). $\tau_a\in\mathcal S$ by the switch lemma ($Y(\tau_a)\le6.21\cdot10^{-13}$, boxes strictly). So $\mathcal S$ is an interval starting at $\tau_a$; let $\tau_e:=\sup\mathcal S$ and let $p$ be the path on $\bigcup_{\tau'\in\mathcal S}[\tau_a,\tau']$.

\emph{(c) Conclusions on $\mathcal S$.} For $\tau'\in\mathcal S$, Steps 2--6 give on $[\tau_a,\tau']$: $t_\tau=\lambda^4(1-\vartheta)$ with $|\vartheta|\le3.37\cdot10^{-8}$; $|x|\le C_xA+x_r+x_{\rm NL}$ (Step 2); $\sqrt{\Lambda_e}\le\bar Y$ (Step 5); the far bounds (Step 6); and $\mathbb B\cup\mathbb B'$ with strict inequalities and the fixed margins of \S``boxes'' (N3$'$ certified at $\bar Y_{\rm box}=1.7\cdot10^{-12}>\bar Y$), with $\bar m\|\hat E\|\le1.04\cdot10^{-4}$. Lemma~\ref{C1:lem:ke} gives $\sup|\tilde\kappa e|\le1.743\cdot10^{-3}$ and $J\ge\frac12$.

\emph{(d) Limits at $\tau_e$.} Suppose $t_e:=\sup_{\tau<\tau_e}t(\tau)<T'$. Since $t_\tau>0$, $t(\tau)\uparrow t_e$. With $a=\frac14+\lambda_\tau/\lambda$ and $b=-c_\tau/\lambda$, Step~1 gives $a=\tilde a+x_1+\vartheta c$, $b=\tilde b-x_2-\vartheta\tilde b$, and by (c) $|x|$ and $\vartheta$ are bounded on $[\tau_a,\tau_e)$, while $\tilde a,\tilde b$ are bounded on the compact label range. So $\ln\lambda$ and $c$ are Lipschitz on $[\tau_a,\tau_e)$ and $p(\tau)\to p_e=(\lambda_e,c_e,t_e)$ with $\lambda_e>0$.

\emph{(e) Closed.} Under the assumption of (d): $\gamma$ is smooth on $[t_{96},t_e]$, so $\Gamma_{p(\tau)}\to\Gamma_{p_e}$ in $C^\infty$, and the base converges with its $x$-jets up to order~4 as $\tau\uparrow\tau_e$. The graph functions $e(\tau)$ obey the uniform box bounds on $|e|,|e'|$ and $|\tilde\kappa e|\le1.743\cdot10^{-3}$, so the foot-point maps and their inverses have derivatives bounded above and below uniformly ($J\ge\frac12$). By Arzel\`a--Ascoli, along a sequence $\tau_n\uparrow\tau_e$, $e(\tau_n)\to e_e$ uniformly and the parameter maps converge to a bi-Lipschitz, increasing, degree-one map $\theta_e$ with $\tilde\Gamma(\cdot,\tau_e)+e_e\tilde\nu(\cdot,\tau_e)=\Gamma_{p_e}\circ\theta_e$. As $|\tilde\kappa e_e|<1$, the normal coordinates of the limit base are a local diffeomorphism along the limit curve, so the foot point is unique and smooth, $\theta_e$ and $e_e$ are as regular as the base, and $e(\tau)\to e_e$ in $C^4$ along the whole limit $\tau\uparrow\tau_e$. The constraints, the box inequalities and $\Lambda_e$, $F_j$ are continuous in this convergence (the boxes and energies involve $x$-jets of $e$ up to order~4 on compact $x$-sets with fixed weights), so they hold at $\tau_e$. Hence $\tau_e\in\mathcal S$, with $p$ extended continuously by $p(\tau_e)=p_e$.

\emph{(f) Open.} If moreover $\tau_e<27.15$: at $(\tau_e,p_e)$, (a) applies with $c_1=1.743\cdot10^{-3}$, and $D_pF(\tau_e,p_e)$ is invertible ($\bar m\|\hat E\|\le1.04\cdot10^{-4}<\frac12$ by (c)). The implicit function theorem for $F$ (continuous in $\tau$, $C^1$ in $p$) gives a continuous extension of $p$ to $[\tau_e,\tau_e+\delta]$ with $F=0$ and $\mathcal G$, $C^1$ on each slab; $t<T'$ for $\delta$ small; and the boxes, strict at $\tau_e$ with fixed margins and continuous in $\tau$, hold on $[\tau_e,\tau_e+\delta]$. Then $\tau_e+\delta\in\mathcal S$, which contradicts $\tau_e=\sup\mathcal S$.

\emph{Conclusion.} If $t_e<T'$, then $\tau_e\in\mathcal S$ by (e) and $\tau_e=27.15$ by (f): alternative (i), with the conclusions of (c) on $[\tau_a,27.15]$. Otherwise $t(\tau)\uparrow T'$ as $\tau\uparrow\tau_e$: alternative (ii), with the conclusions of (c) on $[\tau_a,\tau_e)$. The two alternatives exclude each other ($t(27.15)<T'$ in (i)).
\end{proof}
\begin{remark}[regularity in $\tau$]\label{C1:rem:C1}
Theorem~P2 states that $\lambda,c,t$ are $C^1$. If the base is only continuous in $\tau$ at slab ends (\fn{step4_proof.tex}: ``continuous across slabs''), the implicit function theorem gives continuity at the slab ends and $C^1$ on each slab; that is what Lemma~\ref{C1:lem:P2a} states. Nothing downstream needs more: $t$ is strictly increasing and Lipschitz, $\ln\lambda$ and $c$ are Lipschitz, and the energy inequalities of Steps 4--5 are integrated in $\tau$ (step~4: ``the inequalities hold a.e.\ in $\tau$ and are integrated''). Whether the base is $C^1$ across slab ends was not checked here; flagged for the final review \tH.
\end{remark}

\begin{lemma}[curvature bound on P2's interval]\label{C1:lem:kP2}
In the situation of Lemma~\ref{C1:lem:P2a}, $\sup_{\tau\in[\tau_a,\tau_e)}\int_{\gamma(t(\tau))}k^2ds<\infty$.
\end{lemma}
\begin{proof}\tH, \tS\ Y8. The curve is $ce_1+\lambda(\tilde\Gamma+e\tilde\nu)$ on the half $x\in[0,\tilde\ell]$ and its mirror image. With $x$ the base arclength, $ds_{\rm phys}=\lambda J\,dx$, $k_{\rm phys}=\kappa/\lambda$, and $\kappa=(ae''-e'a'+\tilde\kappa(a^2+e'^2))/J^3$ (Y8b), $a'=-\tilde\kappa'e-\tilde\kappa e'$. Hence $\int k^2_{\rm phys}ds_{\rm phys}=\frac2\lambda\int_0^{\tilde\ell}\kappa^2J\,dx\le\frac{C}{\lambda}\int_0^{\tilde\ell}(1+e''^2)dx$ with $C$ depending on $\sup(|e|,|e'|)$ (boxes), $\sup(|\tilde\kappa|,|\tilde\kappa'|)$ (the base is $C^4$ with enclosed jets on the compact label range) and $J\ge\frac12$ (Lemma~\ref{C1:lem:ke}). (In $x$, with $y=x-X_0$, $X_0=2.2490776515$.) On $x\le31.25$, i.e.\ $y\le29$, where $\chi=1$: $\int_0^{31.25}e''^2\,dx\le\rho(31.25)^{-1}\|(\chi e)''\|_\rho^2\le\rho(31.25)^{-1}\bar Y^2/\varepsilon$ ($\rho$ decreasing in $x$; $\Lambda_e\ge\varepsilon B^2$). On $x\ge30.25$, i.e.\ $y\ge28$, the box gives $x|e''|\le b_1$. These two ranges cover $[0,\tilde\ell]$, so $\int_0^{\tilde\ell}e''^2dx<\infty$ uniformly in $\tau$ (review fix15 C1-3: the v14 text mixed $x$ and $y$; the conclusion is unchanged). $\tilde\ell\le3.7\cdot10^5$. Finally $\lambda_\tau/\lambda=a_{\rm mod}-\frac14$ with $|a_{\rm mod}-\tilde a|\le|x_1|+|\vartheta c|$ bounded (Step~2: $|x|\le C_xA+x_r+x_{\rm NL}$) and $\tilde a$ bounded, so $\lambda\ge\lambda(\tau_a)e^{-C(27.15-\tau_a)}>0$.
\end{proof}

\begin{proposition}[identification on P2]\label{C1:prop:P2}
$T_*>t(27.15)\ge t_1$, and on $[t_{96},t(27.15)]$ the conclusions of Theorem~P2 hold for the smooth flow $\gamma$ of Lemma~\ref{C1:lem:max}.
\end{proposition}
\begin{proof}\tH\ Apply Lemma~\ref{C1:lem:P2a} with $T'=T_*$ (the flow of Lemma~\ref{C1:lem:max} restricted to $[t_{96},T_*)$; $\gamma(t_{96})$ is the P1 curve by Proposition~\ref{C1:prop:P1}). If $t(\tau)\uparrow T_*$ as $\tau\uparrow\tau_e$, then $T_*\le t_{96}+\int_{\tau_a}^{27.15}\lambda^4<\infty$, and $\int k^2ds$ is bounded on $[0,t_{96}]$ (Lemma~\ref{C1:lem:kP1} and Proposition~\ref{C1:prop:P1}) and on $[t_{96},T_*)$ (Lemma~\ref{C1:lem:kP2}), contradicting Lemma~\ref{C1:lem:max}. So $\tau_e=27.15$ and $t(27.15)<T_*$. $t$ is increasing ($t_\tau=\lambda^4(1-\vartheta)>0$), and $t_1=t(\tau^*)$ with $\tau^*\le27.15$ (slab 180).
\end{proof}

%======================================================================
\section{The theorem}
\begin{theorem}[C-1]\label{C1:thm}
Let $\gamma$ be the maximal smooth CDF flow from $\Gamma_0$ given by the Lean theorems (Lemma~\ref{C1:lem:max}). Then $T_*>t(27.15)\ge t_1$ for every good time $\tau^*$ of slab 180, and:
\begin{enumerate}[label=(\roman*),leftmargin=2em]
\item $\gamma$ is jointly $C^\infty$ on $[0,t(27.15)]\times\mathbb S^1$, each $\gamma(t)$ is a closed immersed curve, symmetric: $\overline{\gamma(t)}=\gamma(t)\circ\iota$; this includes $t=0$ and $t=t_1$;
\item on $[0,t_{96}]$, step~4's P1 solution is the tangent angle of $\gamma$ in normalised-arclength gauge, with the same $L$, $X_0$ and clock (Proposition~\ref{C1:prop:P1});
\item on $[t_{96},t(27.15)]$, step~4's P2 decomposition is a decomposition of $\gamma$ (Proposition~\ref{C1:prop:P2});
\item hence every certified hand-off item of Theorem~4(c), (H1)--(H5), is a statement about $\gamma(t_1)$, and (H6) holds: $\gamma(t_1)$ is a smooth closed immersed curve. In the notation of \fn{U1_regularity_v42.tex}, (S4.1) and (S4.2) hold, the latter with $\iota(z)=\bar z$;
\item the maximal flow of \fn{U1_regularity_v42.tex} Lemma ``maximal smooth solution'' (from $F=\gamma(t_1)$) is $t\mapsto\gamma(t_1+t)$, $t\in[0,T_*-t_1)$, i.e.\ its $T_{\max}=T_*-t_1$: step~3 continues the same flow.
\end{enumerate}
\end{theorem}
\begin{proof}\tH\ (i)--(iii) are Lemma~\ref{C1:lem:max} and Propositions~\ref{C1:prop:P1}, \ref{C1:prop:P2} ($\gamma|_{[0,T]}$ is an \fn{IsFlow} for $T<T_*$, in particular on $[0,t(27.15)]$). (iv): step~4's hand-off quantities are functions of the true curve at $t_1$, which by (iii) is $\gamma(t_1)$. (v) \LeanSthree\ (review W7 S3b; text v42): (MS) \fn{cdf_maximal_shift} with $\Phi=\gamma$, $T_{\max}=T_*$ (clauses (i), (ii) of (T1), Lemma~\ref{C1:lem:max}) and $0\le t_1<T_*$ gives: $t\mapsto\gamma(t_1+t)$ is an \fn{IsFlow} on every $[0,T]$, $0<T<T_*-t_1$, every \fn{IsFlow} from $\gamma(t_1)$ on $[0,T]$ has $T<T_*-t_1$ and equals $\gamma(t_1+\cdot)$ there, and the criterion holds if $T_*<\infty$. These are the defining properties of the maximal flow of Lemma ``maximal smooth solution'' of \fn{U1_regularity_v42.tex} (from (T1) for $F=\gamma(t_1)$); by its clause (ii) and the one just stated, the two flows agree and $T_{\max}=T_*-t_1$. (Text v41 cited (L2) and an endpoint gluing lemma with a hand proof.)
\end{proof}
\begin{remark}
This replaces step~4's sentences ``smooth for $t>0$ by parabolic regularity'' (l.672, [A]) and ``continuation uses the literature fact (P1)'' (ll.130, 336, 370, [A]) by Lemma~\ref{C1:lem:max} and Propositions~\ref{C1:prop:P1}, \ref{C1:prop:P2}. It does not use step~4's (H6) or step~3's (P1$'$) as an input, so there is no circularity: (P1$'$) for step~3 is (v) together with \fn{U1_regularity_v42.tex}. Smoothness holds on the closed interval from $t=0$, not only for $t>0$, because $\Gamma_0$ is a trigonometric curve; no parabolic smoothing is used.
\end{remark}

%======================================================================
\section{Gaps, stated plainly}
\begin{enumerate}[leftmargin=*]
\item \textbf{Lean} (text v42): the theorems used are compiled (\fn{U1_regularity_v42.tex} \S1.2; review \fn{writeup/lean_bridge/REVIEW_W7.md}). The two hand bridges of review W7, S3a (the data on \fn{Circle}, Lemma~\ref{C1:lem:max}) and S3b (Theorem~\ref{C1:thm}(v)), are Lean theorems since text v42: (PC), (MS) \LeanSthree. Open: a full \fn{lake build CWFlow.Curve.All}; the commit of \fn{CWFlow/Curve/} and its citation (until then ``commit pending'' with the snapshot SHA256 sums).
\item \textbf{DKS}: no longer an input (text v42). (D1) and the C-4 completion of text v41 are replaced by (T1) (and (T2), which C-1 does not need); the re-reading of DKS is void.
\item \textbf{Step-4 wording} (text v39): the hand-off list and the replaced \textup{[A]} sentences are in \fn{step4_handoff_v42.tex} item (3); \fn{step4_proof.tex} (closed) is not edited, the proposed edits are \fn{STEP4_PATCH_NOTE_v42.md}.
\item \textbf{P2 as an a priori theorem}: Lemma~\ref{C1:lem:P2a} (text v38) writes out the continuity argument \tH; P2's analytic Steps 2--6 keep step~4's status ([D], written out in \fn{LEMMA_F.md}, \fn{T7_REPORT.md}); Remark~\ref{C1:rem:C1} flags the regularity in $\tau$ at slab ends. Lemma~\ref{C1:lem:kP2} uses P2's boxes and Step~2's bound on $|x|$ qualitatively; only the factor $a\ge\frac12$ is a new certificate.
\item \textbf{Galerkin existence} (steps 1--5 of step~4's appendix) is used as stated there (reviewed in step~4). Lemma~\ref{C1:lem:uniq} re-proves only uniqueness.
\item \textbf{Parseval and parity for infinitely many modes} (Lemma~\ref{C1:lem:gauge}(c)) are \tH; the sympy checks Y6, Y7 are on symbolic trigonometric polynomials of degree 3.
\end{enumerate}
\end{document}
